
\documentclass{article}
\usepackage{colortbl}
\usepackage{makecell}
\usepackage{multirow}
\usepackage{supertabular}

\begin{document}

\newcounter{utterance}

\centering \large Interaction Transcript for game `cladder', experiment `full\_v1.5\_default', episode 4386 with Qwen3.5{-}9B.
\vspace{24pt}

{ \footnotesize  \setcounter{utterance}{1}
\setlength{\tabcolsep}{0pt}
\begin{supertabular}{c@{$\;$}|p{.15\linewidth}@{}p{.15\linewidth}p{.15\linewidth}p{.15\linewidth}p{.15\linewidth}p{.15\linewidth}}
    \# & \multicolumn{2}{c}{Player} && \multicolumn{2}{c}{Game Master} \\
    \hline

    \theutterance \stepcounter{utterance}  
    & & & \multicolumn{4}{p{0.6\linewidth}}{
        \cellcolor[rgb]{0.9,0.9,0.9}{
            \makecell[{{p{\linewidth}}}]{
                \texttt{\tiny{[P1$\langle$GM]}}
                \texttt{You are an expert in causal inference. The following question is not a typical commonsense query, but rather a meticulously designed question created by a professor specializing in causal inference, intended to assess the students' mastery of the course content.} \\
\texttt{Imagine a self{-}contained, hypothetical world with only the following conditions, and without any unmentioned factors or causal relationships: Gender has a direct effect on smoking and peanut allergy. Smoking has a direct effect on tar deposit. Tar deposit has a direct effect on peanut allergy. Gender is unobserved. For nonsmokers and with no tar deposit, the probability of being allergic to peanuts is 35\%. For nonsmokers and with high tar deposit, the probability of being allergic to peanuts is 63\%. For smokers and with no tar deposit, the probability of being allergic to peanuts is 25\%. For smokers and with high tar deposit, the probability of being allergic to peanuts is 49\%. For nonsmokers, the probability of high tar deposit is 22\%. For smokers, the probability of high tar deposit is 73\%. For smokers, would it be more likely to see being allergic to peanuts if the person had been a nonsmoker?} \\
\texttt{Start your answer with 'Yes' or 'No', followed by additional reasoning or evidence to support your explanation.} \\
            }
        }
    }
     \\ \\

    \theutterance \stepcounter{utterance}  
    & \multicolumn{4}{p{0.6\linewidth}}{
        \cellcolor[rgb]{0.561,0.737,0.561}{
            \makecell[{{p{\linewidth}}}]{
                \texttt{\tiny{[P1$\rangle$GM]}}
                \texttt{No} \\
\\ 
\texttt{**Reasoning:**} \\
\\ 
\texttt{To determine if a smoker would be more likely to be allergic to peanuts if they had been a nonsmoker, we must calculate the **Counterfactual Probability** \$P(\textbackslash{}text\{Allergy\} \textbackslash{}mid \textbackslash{}text\{Smoker\}, \textbackslash{}text\{Nonsmoker\})\$. This requires integrating over the distribution of the mediator (tar deposit) under the hypothetical condition of being a nonsmoker, while holding the individual's actual characteristics (which include their unobserved gender) constant.} \\
\\ 
\texttt{**1. Define the Structural Relationships and Parameters:**} \\
\texttt{Let \$G\$ be Gender, \$S\$ be Smoking status, \$T\$ be Tar deposit, and \$A\$ be Peanut Allergy.} \\
\texttt{The causal graph is: \$G \textbackslash{}to S\$, \$G \textbackslash{}to A\$, \$S \textbackslash{}to T\$, \$T \textbackslash{}to A\$.} \\
\texttt{We are given the conditional probabilities for \$A\$ given \$S\$ and \$T\$:} \\
\texttt{*   \$P(A \textbackslash{}mid S=0, T=0) = 0.35\$} \\
\texttt{*   \$P(A \textbackslash{}mid S=0, T=1) = 0.63\$} \\
\texttt{*   \$P(A \textbackslash{}mid S=1, T=0) = 0.25\$} \\
\texttt{*   \$P(A \textbackslash{}mid S=1, T=1) = 0.49\$} \\
\\ 
\texttt{We are given the conditional probabilities for \$T\$ given \$S\$:} \\
\texttt{*   \$P(T=1 \textbackslash{}mid S=0) = 0.22\$} \\
\texttt{*   \$P(T=1 \textbackslash{}mid S=1) = 0.73\$} \\
\\ 
\texttt{**2. Analyze the Counterfactual Scenario:**} \\
\texttt{We are evaluating a specific individual who is actually a **Smoker** (\$S=1\$). We want to know the probability of allergy if this person had been a **Nonsmoker** (\$S=0\$).} \\
\texttt{In this counterfactual world, the individual's gender \$G\$ remains unchanged. However, their smoking status changes from 1 to 0. Consequently, their tar deposit distribution changes according to the nonsmoker profile.} \\
\\ 
\texttt{The probability of allergy in this counterfactual scenario is:} \\
\texttt{\$\$P(A \textbackslash{}mid S=1, S\_\{do\}=0) = \textbackslash{}sum\_\{t \textbackslash{}in \textbackslash{}\{0,1\textbackslash{}\}\} P(A \textbackslash{}mid S=0, T=t) \textbackslash{}times P(T=t \textbackslash{}mid S=0)\$\$} \\
\\ 
\texttt{*Note: Since the effect of Gender on Tar Deposit is not explicitly parameterized in the prompt, and the prompt provides the total probability of Tar Deposit given Smoking status (\$P(T|S)\$), we must assume the provided \$P(T|S)\$ values represent the marginal distributions resulting from the specific mix of genders in the population or that the \$T \textbackslash{}to A\$ relationship is sufficient to answer the} \\
            }
        }
    }
    & & \\ \\

    \theutterance \stepcounter{utterance}  
    & & & \multicolumn{2}{p{0.3\linewidth}}{
        \cellcolor[rgb]{0.9,0.9,0.9}{
            \makecell[{{p{\linewidth}}}]{
                \texttt{\tiny{[GM$|$GM]}}
                \texttt{No} \\
            }
        }
    }
    & & \\ \\

    \theutterance \stepcounter{utterance}  
    & & & \multicolumn{2}{p{0.3\linewidth}}{
        \cellcolor[rgb]{0.9,0.9,0.9}{
            \makecell[{{p{\linewidth}}}]{
                \texttt{\tiny{[GM$|$GM]}}
                \texttt{no} \\
            }
        }
    }
    & & \\ \\

    \theutterance \stepcounter{utterance}  
    & & & \multicolumn{2}{p{0.3\linewidth}}{
        \cellcolor[rgb]{0.9,0.9,0.9}{
            \makecell[{{p{\linewidth}}}]{
                \texttt{\tiny{[GM$|$GM]}}
                \texttt{game\_result = WIN} \\
            }
        }
    }
    & & \\ \\

\end{supertabular}
}

\end{document}
